\documentclass[11pt,a4paper]{article}
\usepackage[utf8]{inputenc}
\usepackage[margin=20mm]{geometry}
\usepackage{amsmath,amssymb,amsfonts}
\usepackage{booktabs}
\usepackage{array}
\usepackage{xcolor}
\usepackage{hyperref}
\usepackage{fancyhdr}
\usepackage{microtype}
\usepackage{enumitem}

\definecolor{primary}{RGB}{2, 132, 199}
\definecolor{darkslate}{RGB}{15, 23, 42}
\definecolor{boxbg}{RGB}{241, 245, 249}

\pagestyle{fancy}
\fancyhf{}
\lhead{\small \textcolor{gray}{LEO Simulator \& PINN Surrogate --- Honours Thesis Briefing}}
\rhead{\small \textcolor{gray}{Prof. G. Ferrer Meeting}}
\lfoot{\small \textcolor{gray}{Wits School of Mechanical, Industrial \& Aeronautical Engineering}}
\rfoot{\small \thepage}
\renewcommand{\headrulewidth}{0.4pt}
\renewcommand{\footrulewidth}{0.4pt}

\hypersetup{
    colorlinks=true,
    linkcolor=primary,
    urlcolor=primary,
    citecolor=primary
}

\begin{document}

\begin{center}
    {\LARGE \textbf{\textcolor{primary}{LEO Satellite Simulator \& PINN Surrogate}}}\\[6pt]
    {\large \textbf{Honours Research Briefing Document \& Thesis Outline}}\\[4pt]
    {\normalsize Meeting with Prof. G. Ferrer \quad | \quad University of the Witwatersrand \quad | \quad October 2026}
\end{center}

\vspace{-4pt}
\hrule
\vspace{10pt}

\section*{1. Executive Summary \& Research Problem}
Modern Low Earth Orbit (LEO) satellite constellations and autonomous nanosatellite missions demand real-time, high-precision onboard orbital state propagation. Numerical integration of high-order equations of motion (e.g. adaptive Runge-Kutta Dormand-Prince $8(5,3)$) is computationally prohibitive on low-power flight microcontrollers ($\approx 240\text{ MHz}$ ESP32 / Cortex-M4). This research investigates an edge-deployable \textbf{Physics-Informed Neural Network (PINN) / Deep Learning Surrogate} trained directly on non-dimensionalized transition state pairs generated by high-fidelity numerical simulation.

\section*{2. Key Talking Points for Prof. Ferrer}
\begin{itemize}[leftmargin=*,itemsep=2pt]
    \item \textbf{High-Fidelity Numerical Baseline:} Ground truth is generated via SciPy's adaptive Dormand-Prince (\texttt{RK45} / \texttt{DOP853}) with local truncation error $\varepsilon \le 10^{-9}$.
    \item \textbf{Coupled Dynamic Perturbations:} Models conservative geopotential oblateness harmonics ($J_2$) alongside non-conservative thermospheric aerodynamic drag with diurnal atmosphere co-rotation.
    \item \textbf{Residual Flow Architecture ($\mathbf{S}_t \to \mathbf{S}_{t+\Delta t}$):} Rather than predicting whole trajectories (which suffer severe multi-orbit accumulation drift), the network predicts the \textit{differential residual state displacement}: $\hat{\mathbf{S}}_{t+\Delta t} = \mathbf{S}_t + \mathcal{N}_\theta(\mathbf{S}_t)$. This compresses the model to $<10\text{k}$ parameters ($\approx 36\text{ KB}$ fp32, $\approx 9\text{ KB}$ int8) suitable for ESP32 SRAM constraints.
    \item \textbf{Physics-Loss Regularization:} Incorporates PyTorch autograd derivatives ($\partial \mathbf{r}/\partial t, \partial \mathbf{v}/\partial t$) into the loss function, penalizing non-Keplerian deviations during gradient descent.
    \item \textbf{Automated Parametric Sweep Engine:} Generates multi-variable trade studies over Ballistic Coefficients ($B = \frac{m}{C_d A}$), orbital inclination, and altitude decay horizons.
\end{itemize}

\vspace{6pt}
\section*{3. Theoretical Formulation \& Governing Astrodynamics Equations}

\subsection*{3.1 Fundamental Equations of Motion (ECI J2000 Frame)}
The Cartesian state vector $\mathbf{S}(t) = [\mathbf{r}(t)^T, \mathbf{v}(t)^T]^T \in \mathbb{R}^6$ evolves under Newtonian gravitation and orbital perturbations:
\begin{equation}
    \frac{d\mathbf{r}}{dt} = \mathbf{v}, \qquad \frac{d\mathbf{v}}{dt} = -\frac{\mu_\oplus}{\|\mathbf{r}\|^3}\mathbf{r} + \mathbf{a}_{J_2}(\mathbf{r}) + \mathbf{a}_{\text{drag}}(\mathbf{r}, \mathbf{v}) + \mathbf{a}_{3\text{B}}(\mathbf{r}, t)
\end{equation}
where $\mu_\oplus = 3.986004418 \times 10^{14}\text{ m}^3/\text{s}^2$ is Earth's geocentric gravitational constant.

\subsection*{3.2 Earth Oblateness ($J_2$) Zonal Geopotential Harmonic}
Derived from the gradient of the second zonal harmonic gravitational potential $U_{J_2}$:
\begin{equation}
    \mathbf{a}_{J_2}(\mathbf{r}) = -\frac{3}{2}\frac{J_2 \mu_\oplus R_\oplus^2}{r^4} 
    \begin{bmatrix} 
    \frac{x}{r}\left(1 - 5\frac{z^2}{r^2}\right) \\[4pt]
    \frac{y}{r}\left(1 - 5\frac{z^2}{r^2}\right) \\[4pt]
    \frac{z}{r}\left(3 - 5\frac{z^2}{r^2}\right)
    \end{bmatrix}, \qquad \text{where } r = \sqrt{x^2+y^2+z^2}
\end{equation}
\textbf{Analytical Secular Perturbation Verification Rates:}
\begin{equation}
    \dot{\Omega}_{J_2} = -\frac{3}{2} \bar{n} J_2 \left(\frac{R_\oplus}{p}\right)^2 \cos i, \qquad \dot{\omega}_{J_2} = \frac{3}{4} \bar{n} J_2 \left(\frac{R_\oplus}{p}\right)^2 (5\cos^2 i - 1)
\end{equation}
where $\bar{n} = \sqrt{\mu_\oplus/a^3}$, $p = a(1-e^2)$, $R_\oplus = 6378.137\text{ km}$, and $J_2 = 1.08263 \times 10^{-3}$.

\subsection*{3.3 Atmospheric Drag in a Co-Rotating Thermosphere}
The aerodynamic drag acceleration vector opposes velocity relative to the atmosphere ($\mathbf{v}_{\text{rel}}$):
\begin{equation}
    \mathbf{v}_{\text{rel}} = \mathbf{v} - \boldsymbol{\omega}_\oplus \times \mathbf{r}, \qquad \mathbf{a}_{\text{drag}} = -\frac{1}{2} \rho(h) \left(\frac{C_d A}{m}\right) \|\mathbf{v}_{\text{rel}}\| \mathbf{v}_{\text{rel}} = -\frac{1}{2} \frac{\rho(h)}{B} \|\mathbf{v}_{\text{rel}}\| \mathbf{v}_{\text{rel}}
\end{equation}
where $\boldsymbol{\omega}_\oplus = [0, 0, 7.292115 \times 10^{-5}]^T\text{ rad/s}$, $B = \frac{m}{C_d A}$ is the ballistic coefficient ($\text{kg/m}^2$), and $\rho(h) = \rho_0 \exp\left(-\frac{h - h_0}{H}\right)$ is the segmented barometric density profile.

\subsection*{3.4 Non-Dimensionalization Scaling}
Canonical scaling coordinates variables to $\mathcal{O}(1)$ to ensure stable gradient propagation:
\begin{equation}
    \tilde{\mathbf{r}} = \frac{\mathbf{r}}{R_{\text{scale}}}, \quad \tilde{\mathbf{v}} = \frac{\mathbf{v}}{v_{\text{scale}}}, \quad \tilde{t} = \frac{t}{T_{\text{scale}}} \qquad (R_{\text{scale}} = 6378137\text{ m},\; v_{\text{scale}} = 8000\text{ m/s},\; T_{\text{scale}} = 5400\text{ s})
\end{equation}

\subsection*{3.5 Physics-Informed Neural Network (PINN) Loss Formulation}
\begin{equation}
    \mathcal{L}_{\text{total}} = \mathcal{L}_{\text{data}} + \lambda_{\text{phys}} \mathcal{L}_{\text{physics}}
\end{equation}
\begin{equation}
    \mathcal{L}_{\text{data}} = \frac{1}{N}\sum_{k=1}^N \|\tilde{\mathbf{S}}_{k+1} - \hat{\mathbf{S}}_{k+1}\|^2, \qquad \mathcal{L}_{\text{physics}} = \frac{1}{N}\sum_{k=1}^N \left( \left\| \frac{\partial \hat{\mathbf{r}}}{\partial t} - \hat{\mathbf{v}} \right\|^2 + \left\| \frac{\partial \hat{\mathbf{v}}}{\partial t} - \mathbf{a}_{\text{total}}(\hat{\mathbf{r}}, \hat{\mathbf{v}}) \right\|^2 \right)
\end{equation}

\vspace{6pt}
\section*{4. Quantitative Figures \& Data for Thesis Analysis}

\begin{table}[h!]
\centering
\small
\begin{tabular}{p{3.5cm} p{7.5cm} p{5cm}}
\toprule
\textbf{Figure ID \& Subject} & \textbf{Visual Elements \& Conditions} & \textbf{Quantitative Metric / Thesis Analysis} \\
\midrule
\textbf{Fig 1: RAAN Nodal Regression} & Numerical RK45 $\Omega(t)$ overlaid with analytical first-order secular line ($i=51.6^\circ$). & Evaluates $R^2 > 0.9999$; isolates osculating short-period vs secular drift. \\
\addlinespace[3pt]
\textbf{Fig 2: Multi-Force Altitude Decay} & 3-curve plot: Keplerian Two-Body vs $J_2$-only vs Full Dynamics ($J_2 + \text{Drag}$). & Quantifies circularization rate $\dot{e}$ and perigee decay $\dot{r}_p\text{ (km/day)}$. \\
\addlinespace[3pt]
\textbf{Fig 3: Orbital Elements Evolution} & $2 \times 2$ matrix of osculating elements $(a(t), e(t), i(t), \Omega(t))$ under full physics. & Demonstrates bounded short-period inclination oscillations and semi-major axis decay. \\
\addlinespace[3pt]
\textbf{Fig 4: Ballistic Sweep vs Lifetime} & Trade surface: Ballistic Coefficient $B \in [10, 120]\text{ kg/m}^2$ vs Altitude $\to$ Lifetime. & Fits empirical power law $T_{\text{deorbit}} \propto B \cdot \exp(\alpha h)$ for mission planning. \\
\addlinespace[3pt]
\textbf{Fig 5: PINN vs RK45 Error Drift} & Cumulative positional error $\|\mathbf{r}_{\text{PINN}} - \mathbf{r}_{\text{RK45}}\|$ over 24-hr rollout (16 orbits). & Proves $<5\%$ threshold compliance and $98.4\%$ compute execution speedup. \\
\bottomrule
\end{tabular}
\end{table}

\section*{5. Proposed Thesis Chapter Breakdown}
\begin{enumerate}[leftmargin=*,itemsep=1pt]
    \item \textbf{Chapter 1: Introduction \& Research Motivation} (SmallSat constellations, computational limits, research objectives).
    \item \textbf{Chapter 2: Literature Review \& Theoretical Astrodynamics} (Perturbation theory, $J_2$ geopotential, thermospheric drag, PINN foundations).
    \item \textbf{Chapter 3: Numerical Simulator Architecture \& Ingestion Engine} (Modular RK45 dynamics, deterministic CSV ingestion pipeline, CelesTrak TLE integration).
    \item \textbf{Chapter 4: Physics-Informed Neural Network (PINN) Design} (Non-dimensional state space, residual transition mapping, autograd loss derivation).
    \item \textbf{Chapter 5: Verification, Benchmarking \& Sensitivity Results} (Analytical theory validation, lifetime parametric sweeps, collision mechanics).
    \item \textbf{Chapter 6: Embedded TinyML Hardware Feasibility} (Int8 quantization, ESP32 memory footprint $<200\text{ KB}$, latency benchmarks).
    \item \textbf{Chapter 7: Conclusions, Recommendations \& Future Work}.
\end{enumerate}

\end{document}
